\cotefile{91}
\scanpage{87}

\margin{[sous-catégorie pleine de (\unsure{Sch|Schs})]}
\diagram{\scriptsize
\begin{tikzcd}[ampersand replacement=\&,row sep=large]
{\text{Préschémas atomiques : (sommes de schémas de corps)}} \arrow[d, "{\substack{\text{foncteur \unsure{pleinement|pleinmt} fidèle,}\\\text{commutant aux sommes}}}"]\\
{\substack{\text{Foncteurs covariants}\quad(\text{Corps})\unsure{\longrightarrow|\text{---}}(\text{Ens})\\
\text{i.e. contravariants}\quad(\text{Schs ponctuels})^{o}\longrightarrow(\text{Ens})}}
\end{tikzcd}}

\diagram{\scriptsize
\begin{tikzcd}[ampersand replacement=\&,row sep=large]
{\substack{\text{Préschémas atomiques sur } k\quad [\text{i.e. sur }\operatorname{Spec}(k)]\quad
\deleted{(At)}/\operatorname{Spec}k}}
\arrow[d,"{\text{foncteur \unsure{pleinement|pleinmt} fidèle}}"]\\
{\substack{\text{Foncteurs covariants}\quad(\text{Extensions de }k)\longrightarrow(\text{Ens})\\
\text{Foncteurs ''}\qquad(\text{Schs ponctuels}/\operatorname{Spec}k)^{o}\longrightarrow(\text{Ens})}}
\end{tikzcd}}

Soit $X$ un préschéma, il définit un préschéma atomique $|X|=\coprod_{x\in X}\operatorname{Spec}k(x)$, qui \deleted{\unsure{est représenté|est égal}}\inserted{\unsure{correspond|représente}} \unsure{à|:} la restriction du foncteur contravariant défini par $X:(\text{Sch})^{o}\longrightarrow(\text{Ens})$, aux schémas ponctuels.

Bien sur, on obtient un foncteur $X\leadsto |X|$, dont la \uline{restriction aux préschémas réduits} \uline{est} \uline{fidèle}.

\deleted{\unsure{principes|principaux}}

Soient $X,Y$ deux préschémas, on appelle \uline{pseudo-morphisme} ou \uline{application constructible} de $X$ dans $Y$, tt morphisme $f:|\unsure{\overline X|X}|\longrightarrow|\unsure{\overline Y|Y}|$ qui a la propriété suivante : Pour tte partie \deleted{quasi-compact}\inserted{ouverte affine} \deleted{\unsure{donné|affine|de}} $U$ de $X$, il existe une partition de $U$ en des $U_i$ \inserted{\uline{loc. fermés}} \uline{constructibles}, tels que $f\mid |U_i|$ soit un \uline{morphisme} ${U_i}_{\text{réd}}\to Y$.
% archive sheet number: 86

\scanpage{88}

\mspage{- 139 -}\hfill n\textsuperscript{o} 316

a) Soit $M'_s$ le sous-espace de $M_s$ engendré par les éléments de la forme $\beta_1\otimes\beta_2\otimes\ldots\otimes\beta_s$, où $\beta_i\in\mathcal{L}(V)$ et où $\beta_i=1$ pour au moins un $i$. Soient $z$ un élément de $U_s$, $z_{\unsure{o|0}}=\sum_{1\leq i_1,\ldots,i_s\leq m}a_{i_1..i_s}\alpha_{i_1}\ldots\alpha_{i_s}$ sa composante dans $U^s$, où les $a_{i_1\ldots i_s}$ sont symétriques par rapport aux permutations des indices. Montrer que
\[
\begin{aligned}
&\pi_s(z)\equiv s!\\
&\sum_{1\leq i_1\ldots,i_s\leq m}a_{i_1\ldots i_s}\,\alpha_{i_1}\otimes\ldots\otimes\alpha_{i_s}\quad(\mathrm{mod.}\ M'_s)\,.
\end{aligned}
\]

b) Déduire de a), que, pour tout $z\in U$, il existe une représentation de dimension finie $\pi$ de $U$ telle que $\pi(z)\ne 0$. (Pour une suite de cet exerc., cf. § 7, exerc. 3).

6) - On suppose que $K$ est un corps. Soient $g$ une $K$-algèbre de Lie, $x\longrightarrow x_M$ une représentation de $g$ de dimension finie, $(x_1,\ldots,x_n)$ une base de $M$; et $(x^*_1,\ldots,x^*_n)$ la base duale.

a) Si $f$ est une forme linéaire sur $\overset{r}{\otimes}\,M$, on a :
\[
f=\sum_{1\leq i_1,\ldots,i_r\leq n}f(x_{i_1}\otimes\ldots\otimes x_{i_r})(x^*_{i_1}\otimes\ldots\otimes x^*_{i_r})\,.
\]

b) Soit $\beta$ une forme bilinéaire non dégénérée sur $M$, invariante pour $g$. Soit $(x'_1,\ldots,x'_n)$ la base de $M$ telle que $\beta(x'_i,x'_j)=\delta_{ij}$. Soit $f$ une forme linéaire invariante sur $\overset{r}{\otimes}\,M$. Déduire de a) que
\[
\sum_{1\leq i_1,\ldots,i_r\leq n}f(x_{i_1}\otimes\ldots\otimes x_{i_r})(x'_{i_1}\otimes\ldots\otimes x'_{i_r})
\]
est un élément invariant de $\overset{r}{\otimes}\,M$ indépendant du choix de la base $(x_i)$.

c) Soit $U$ l'algèbre enveloppante de $g$., $U^*$ l'espace dual de l'espace $U$. La représentation adjointe de $g$ se prolonge en une représentation $x\longrightarrow x_U$ de $g$ dans $U$ définie par $x_Uu=xu-ux$ pour tout $u\in U$. Donc $U^*$ est muni d'une structure de $g$-module. Pour qu'un élément $f$ de $U^*$ soit invariant, il faut et il suffit que $f(uv)=f(vu)$ quels que soient $u,v\in U$.
% archive sheet number: 87

\scanpage{89}

\margin{Transitivité $\swarrow$}

\deleted{[N.B. il suffit de le vérifier pour des ouverts \deleted{\unsure{affi}} affines quasi compacts \inserted{$U^\alpha$} recouvrant $X$, et pour tt $U^\alpha$ trouver un \uline{recouvrement} de $U^\alpha$ par des parties \deleted{constructibles}\inserted{loc. fermées} $U_i^\alpha$, tels que $f\mid |U_i^\alpha|$ soit \deleted{\unsure{aff}} morphique \ldots}

\margin{\deleted{Proposition}}

On appelle \uline{ensemble} \deleted{\unsure{algébrique}}\deleted{\unsure{arithmétique}} \uline{schématique} un \deleted{préschéma}\inserted{préschéma quasi-séparé quasi-compact}, considéré comme objet d'une catégorie dont les morphismes sont les applications constructibles.

\diagram{\tiny
\begin{tikzcd}[ampersand replacement=\&,column sep=6pt,row sep=1pt,
execute at end picture={\draw[thin] ([yshift=-8pt]\tikzcdmatrixname-1-1.south east) -- ++(6pt,5pt) -- ++(-4pt,-9pt) -- ([yshift=-6pt]\tikzcdmatrixname-1-2.south west) -- ++(-4pt,-9pt) -- ([yshift=-6pt]\tikzcdmatrixname-1-2.south) -- ++(-3pt,-10pt) -- ([yshift=-6pt]\tikzcdmatrixname-1-3.south west) -- ++(-3pt,-10pt) -- ([yshift=-9pt]\tikzcdmatrixname-1-4.south west);}]
{(\text{Schs }\deleted{\unsure{\text{réd}}})}\arrow[r]
\arrow[rr,bend right=23,"{\deleted{\text{fidèle}}}"']
\arrow[rrr,bend right=30,"{\deleted{\text{fidèle}}}"']
\& {(\text{Schs})}\arrow[r,"{\unsure{\text{n'd}|\text{réd}|\text{red}}\ \text{fidèle}}"']
\arrow[rr,bend right=21,"{\deleted{\text{fidèle}}}"']
\& {(\text{ens Schs})}\arrow[r,"{\text{fidèle}}"']
\& {(\text{Sch. \unsure{atomes|atoms|atom.}})}\arrow[r,"{\text{pl. fid.}}"']
\& {\text{\uline{Hom}}\,((\text{Corps}),\text{Ens})}
\end{tikzcd}}

\margin{\fbox{\parbox{1.8cm}{N.B. Tous les foncteurs écrits \unsure{ici|ci-dessus} commutent aux $\underset{\leftarrow}{\lim}$ finies}}}

On s'intéresse \unsure{donc|données} aux notions sur les ens. schématiques, qui sont invariantes par \unsure{isom|isom.} d'ens schématiques (i.e. isom. constructibles).

Notons \deleted{\unsure{y|$X$}} comme \unsure{Rem. que|lemme que}

\diagram{\begin{tikzcd}[ampersand replacement=\&,row sep=2pt]
X\arrow[r]\&Y\\
X_i\&Y_i\\
X'_j\&\\
X_i\cap X'_j\arrow[r]\&Y_{ij}
\end{tikzcd}}

\uline{Lemme} Soit $X$ \deleted{\unsure{pré}}préschéma quasi-compact, quasi-séparé. \deleted{Alors il existe une} $(Z_i)$ une famille \deleted{\unsure{de|finie|disjointe}} de parties constructibles recouvrant $X$. Alors il existe une partition de $X$ en des parties loc. fermées constructibles affines, plus fine que \unsure{un|une} \unsure{rec. fini|recouvrement} \uline{\unsure{donné}} de $X$ par des ens. constructibles

\uline{Proposition} Supposons $X$ quasi-séparé et quasi-compact, $f:|X|\to|Y|$ donné. Conditions équivalentes

a) $f$ est constructible

\unsure{aiii)|a ii)|(ii)} $\exists$ une partition \inserted{finie} de $X$ en des parties loc. \uline{fermées constructibles} \inserted{affines} $Z_i$, telles que $f\mid |Z_i|$ soit morphique

\deleted{\unsure{(iii)|bii)} Il existe une suite croissante d'ouverts quasi-compacts $U_i$ couvrant $X$, telle que $f\mid |U_i-U_{i-1}|$ soit \unsure{morphique}.}
% archive sheet number: 88

\scanpage{90}

\mspage{- 138 -}\hfill n\textsuperscript{o} 316

dérivation de $V$, alors $s$ et $n$ sont des dérivations de $V$.

4) - On suppose que $K$ est un corps. Soit $g$ une $K$-algèbre de Lie.

a) Soient $x\longrightarrow x_M$ et $x\longrightarrow x_N$ des représentations de $g$, $x\longrightarrow x_N$ étant irréductible et de dimension finie. Soit $(f_j)_{1\leq j\leq n}$ une base de $N$. Si des éléments non tous nuls $e_1,\ldots,e_n$ de $M$ sont tels que $\sum_{i=1}^{n}e_i\otimes f_i$ soit invariant dans $M\otimes N$, alors le sous-espace $M_1$ de $M$ engendré par les $e_i$ est stable pour les $x_M$, et la représentation $x\longrightarrow x_{M_1}$ est duale de la représentation $x\longrightarrow x_N$.

b) Soient $x\longrightarrow x_{M_1}$, $x\longrightarrow x_{M_2}$ deux représentations irréductibles de dimension finie de $g$. Soient $x\to x_N$, $x\to x_{\unsure{N^*|N*}}$ deux représentations duales de dimension finie\unsure{,|{}} de $g$. Pour que le $g$-module $M_1$ soit isomorphe à un sous-$g$-module de $M_2\otimes N$, il faut et il suffit que $M_2$ soit isomorphe à un sous-$g$-module de $M_1\otimes N^*$(Utiliser a)).

c) Si $M_1$ et $M_2$ sont des $g$-modules simples non isomorphes, et si $K_1$ est une extension de $K$, il n'existe pas de $g_{(K_1)}$-module simple isomorphe à la fois à un sous-$g_{(K_1)}$-module de $M_{1(K_1)}$ et à un sous-$g_{(K_1)}$-module de $M_{2(K_1)}$. (Utiliser a)).

\margin{\handwritten{$\unsure{\pi|\Pi}$}}
5) - On suppose que $K$ est un corps de caractéristique 0. Soient $V$ un $K$-espace vectoriel de dimension finie, $g=\mathrm{sl}(V)$, $U$ l'algèbre enveloppante de $g$, $U_n$ l'ensemble des éléments de filtration $\leq n$ dans $U$, $U^n$ l'ensemble des éléments symétriques homogènes d'ordre $n$ dans $U$ $(\alpha_1,\ldots,\alpha_m)$ une base de l'espace vectoriel $g$. Soit $W_s=\overset{s}{\otimes}\,V$. La représentation identique de $g$ définit une représentation de $g$ dans $W_s$ qui se prolonge en un homomorphisme $\pi_s$ de $U$ dans l'algèbre $M_s=\mathcal{L}(W_s)$.
% archive sheet number: 89

\scanpage{91}

\uline{Proposition} Soit $f:X\unsure{\text{---}|\to}Y$ \deleted{morphisme} application constructible d'ens. schématiques. Pour tte partie constructible $E$ de $Y$, $f^{-1}(E)$ est constructible.

\diagram{\begin{tikzcd}[ampersand replacement=\&,row sep=tiny,execute at end picture={\draw[thin] (\tikzcdmatrixname-1-2.west) -- (\tikzcdmatrixname-2-1.center) -- (\tikzcdmatrixname-1-1.east); \draw[thin] (\tikzcdmatrixname-2-1.west) -- (\tikzcdmatrixname-1-1.east);}]
\deleted{X}\arrow[r,"{\deleted{\unsure{f|{}}}}"]\&\deleted{Y}\\
\deleted{X_i}\arrow[u,phantom,"{\deleted{\cup}}" description]\&
\end{tikzcd}}

En effet, soit $X=\coprod_i X_i$, les $X_i$ loc. fermés et constructibles dans $X$, \deleted{\unsure{donc|tels que}} $f_i=(f\mid |X_i|)$ un morphisme. On a $f^{-1}(E)=\bigcup f_i^{-1}(E\unsure{f|i|{}})$, il suffit de prouver que chaque $f^{-1}(E)=f^{-1}(E)\cap X_i$ est constructible dans $\alpha$, ou -- ce qui revient au $\overline{\mathrm m}$, dans $X_i$. [\deleted{\unsure{par}} \unsure{faire|par} un th. de Chevalley, \inserted{\unsure{l'immers.}} $X_i\unsure{\text{---}|\to}X$ étant de présentation finie, pour une structure induite convenable sur $X_i$ \ldots] Cela nous ramène au cas d'un \uline{morphisme} \inserted{\uline{de préschémas}} ; comme il est quasi-compact quasi-séparé, c'est connu.

% slanted
\margin{\unsure{fragile|fragilité|toujours} \unsure{pour|point} l'argument!}

\uline{Corollaire} Invariance de la notion de partie constructible d'un ens schématique

\uline{Proposition} Soit $f:X\longrightarrow Y$ une application constructible de \deleted{ens. schématiques}\inserted{préschémas quasi compacts quasi-séparés}. Pour que $f$ soit un isom d'ens schém., il f. et s. que il existe \deleted{\unsure{de|pour}} un \deleted{\unsure{p}} ens fini $I$, deux partitions $(X_i)_{i\in I}$ et $(Y_j)_{j\in I}$ de $X$ ($Y$) en parties loc. fermées constructibles, \inserted{[affines si on veut]}, telles que $f$ induise un \uline{isom} \uline{de} $X_i$ \uline{sur} $Y_i$.

C'est \unsure{évidem\textsuperscript t|évident} suffisant, \unsure{prouvons|pour} la nécessité. Soit $g$ l'inverse de $f$, $X_\alpha,Y_\beta$ des partitions \inserted{de $X,Y$} mettant en évidence le caractère constructible de $f,g$. \unsure{Évidem\textsuperscript t|Evidemment} $f,g$ sont des bijections inv. l'une de l'autre

\diagram{\begin{tikzcd}[ampersand replacement=\&]
X\arrow[r,shift left,"f"]\&Y\arrow[l,shift left,"g"]
\end{tikzcd}}
% archive sheet number: 90

\scanpage{92}

\mspage{- 137 -}\hfill n\textsuperscript{o} 316

- \uline{Exercices \unsure{srle|surle} §3} -

1) - Soient $g$ une algèbre de Lie, $\unsure{\varrho|\rho}$ une représentation de $g$ dans un $K$-module $M$, $\sigma$ la représentation associée de $g$ dans l'algèbre tensorielle de $M$. Montrer que le sous-module des tenseurs symétriques et le sous-module des tenseurs antisymétriques sont stables pour $\sigma$.

2) - Soient $g$ une algèbre de Lie, $\varrho$ une représentation de $g$ dans un $K$-module $M$, $\sigma$ la représentation associée de $g$ dans $Q=\mathcal{L}(M,M;M)$. Pour qu'un $f\in Q$ soit invariant pour $\sigma$, il faut et il suffit que les $x_M$ soient des dérivations de $M$ muni de la multiplication définie par $f$.

3) - Soit $V$ un espace vectoriel de dimension finie sur un corps parfait $K$.

a) La représentation identique de $\mathrm{gl}(V)$ définit canoniquement des repré\unsure{senta|ssnta|snta}tions de $\mathrm{gl}(V)$ dans $\overset{p}{\otimes}\,V$ et $\overset{q}{\otimes}\,V^*$, donc une représentation $u\longrightarrow u^p_q$ de $\mathrm{gl}(V)$ dans $V^p_q=(\overset{p}{\otimes}\,V)\otimes(\overset{q}{\otimes}\,V^*)$. Montrer que $u^p_q$ est semi-simple si $u$ l'est. (Se ramener, par extension du corps de base, au cas où $K$ est algébriquement clos). Montrer que $u^p_q$ est nilpotent si $u$ l'est. En déduire que, si $s$ et $n$ sont les composantes semi-simple et nilpotente de $u$, $s^p_q$ et $n^p_q$ sont les composantes semi-simple et nilpotente de $u^p_q$.

b) Déduire de a) que, si un élément de $V^p_q$ est annulé par $u^p_q$, il est annulé par $s^p_q$ et $n^p_q$.

c) Déduire de b) et de l'exerc.2 que, si $V$ est muni d'une structure d'algèbre non nécessairement associative, et si $u$ est une
% archive sheet number: 91

\scanpage{93}

Soit $X_{\alpha\beta}=X_\alpha\cap f^{-1}(Y_\beta)$, $Y_{\deleted{\unsure{\alpha\beta}}\beta\alpha}=Y_\beta\cap g^{-1}(X_\alpha)$, donc $X_{\alpha\beta}$ et $Y_{\beta\alpha}$ sont des ens \uline{constructibles}, qui se correspondent par $f,g$.

\deleted{Donc \inserted{\unsure{disons}} $X_{\alpha\beta}$ est une union finie $\bigcup_\lambda X_{\alpha\beta\lambda}$, où $X_{\alpha\beta\lambda}$ est un sous-schéma affine constructible de $X_{\alpha\beta}$ \inserted{\unsure{donc de} $X_\alpha$}. Alors $g^{-1}(X_{\alpha\beta\lambda})=g_\beta^{-1}(X_{\alpha\beta\lambda})$ est \deleted{\unsure{une partie|un sous schéma}} un sous-schéma de $Y_\beta$ constructible.}

\diagram{\begin{tikzcd}[ampersand replacement=\&,row sep=2pt,column sep=small]
X\&Y\\
{}\&Y_\beta\arrow[ul]
\end{tikzcd}}

D'ailleurs $X_{\alpha\beta}=f_\alpha^{-1}(Y_\beta)$ est un \deleted{\unsure{sous-schéma de|ens. schématique de}} loc. fermé dans $X_\alpha$, donc dans $X$. et $f_{\alpha\beta}:X_{\alpha\beta}\to Y$ \deleted{\unsure{est}}est un morphisme \inserted{de schémas} (car induit par $f_\alpha:X_\alpha\to Y$]. \deleted{\unsure{comme il en est}}\inserted{De $=$,} les $Y_{\alpha\beta}$ sont constructibles et loc. fermés, et \inserted{$g_{\alpha\beta}:Y_{\alpha\beta}\to X$ est un morphisme de schém.} \deleted{\unsure{comme}}\inserted{Enfin} $X_{\alpha\beta}\unsure{\text{---}|\to}Y_{\unsure{\alpha|\delta|\beta}}$ se factorise par un \uline{morph.} $X_{\alpha\beta}\to Y_{\alpha\beta}$ [car il le fait \unsure{ens|ainsi}], et de $=$ on a un \uline{morphisme} $Y_{\alpha\beta}\to X_{\alpha\beta}$. Ces derniers sont donc inverses l'un de l'autre.

\uline{Cor.} 0 Si $X$ noethérien $Y$ l'est
% slanted
\margin{\uline{Notion} d'ens schématique \uline{noethérien} [défini par un $X$ noethérien, \unsure{ou|qui} \illegible{1} \unsure{qu'on peut} \unsure{supposer} \unsure{affine}]}

\uline{Cor.} 1 Tout ens schématique est isomorphe à un $|X|$, avec $X$ affine.

\uline{Cor.} 2 Tout morphisme d'ens\unsure{.} schématiques est isom. à un morphisme de \deleted{ens} schémas affines. (\unsure{Considéré|Considérés} comme application \unsure{constructible|constructibles}).

Démonstration \unsure{et|tte} \unsure{réciproque|analogue|intrinsèque|logiques}.

\deleted{\uline{Cor.} 3 Un \deleted{isom}\inserted{morphisme $f:X\to Y$} de préschémas quasi compacts et quasi-séparés, qui est un isom d'ens schématiques est \uline{de type fini}}
% archive sheet number: 92

\scanpage{94}

\mspage{- 37 -}\hfill n\textsuperscript{o} 288

dont les projections sur les facteurs $E_i$ sont bornées, que les parties bornées de $E_i$ sont contenues (du moins si $K$ est non discret) dans un homothétique de la boule unité de $E_i$, il s'ensuit aisément que sur l'espace vectoriel $\mathrm{L}(E_1,\ldots,E_n\,;\,F)$, la topologie de la convergence bornée est définie par l'unique norme :
\[
\lVert u\rVert=\operatorname{Sup}\ \lVert u(x_1,\ldots,x_n)\rVert\quad\text{pour}\quad\lVert x_1\rVert<1,\ldots,\lVert x_n\rVert\leq 1
\]

Ceci mérite sans doute une proposition en forme et une rédaction moins cursive, mais l'endroit pour ces sorites semblerait plutôt dans les E.V.T.
% archive sheet number: 93

\scanpage{95}

\uline{Morphisme de type fini} \deleted{[\unsure{prés. finie|quasi-fini}]} d'ens schématiques : tel qu'il existe une partition de $X$ en ensembles loc. fermés constructibles affines $X_i$, telle que les $f\mid X_i$ soient des morphismes de type fini \deleted{[\unsure{prés. finie|quasi-finis}]}

% slanted
\margin{est-il indiqué de définir la présentation finie ?}

Un tel morphisme de \unsure{'ens|l'ens|un ens.} schématique est isom \unsure{:|à} un morphisme de type fini de schémas affines, Considéré comme morphisme d'ens schématiques). \uline{Transitivité}.

Ens schématique sur un autre, de type fini sur un autre. Tous sont sur $\mathbb{Z}$ !

Ens schématique défini par une partie constructible d'un ens. schématique, son immersion canonique, décomposition de $X$ en somme $X=E\amalg(X-E)$, critère de factorisabilité par une partie constructible ; c'est \uline{ensembliste} !

Produit\unsure{s} fibrés [et $\underset{\leftarrow}{\lim}$ finies] dans la catégorie des ens schématiques. \uline{Existe}, et le foncteur \deleted{\unsure{préschémas}} $(\text{présch }\unsure{\text{q. compact}|\text{qu compact}}\ \text{q. sép})\longrightarrow(\text{ens schém})$ commute aux produits fibrés.

Si on a $X$ et $Y$ sur $Z$, si $X',Y'$ sont des parties constructibles de $X,Y$ alors $X'\times_Z Y'=\mathrm{pr}_1^{-1}(X')\cap\mathrm{pr}_2^{-1}(Y')$ \ldots.
% archive sheet number: 94

\scanpage{96}

\hfill n\textsuperscript{o} 285

\begin{center}
RAPPORT sur la R\handwritten{\'}{E}\handwritten{\'}{E}DITION\\
de\\
TOPOLOGIE GENERALE\\
Chap.III
\end{center}

§2 ; les n\textsuperscript{os} 4 et 8 sont remaniès de la façon suivante :

4. - \uline{Espaces homogènes}

Soient $E$ un espace topologique, $G$ un groupe topologique, et supposons $E$ muni du groupe d'opérateurs $G$ : (Alg. chap.I, § 7, n\textsuperscript{o}2) ; nous dirons que $G$ \uline{opère continûment} dans $E$ si l'application $(s,x)\to s.x$ de $G\times E$ dans $E$ est \uline{continue}. Lorsqu'il en est ainsi, pour tout $s\in G$, l'application $x\longrightarrow s.x$ de $E$ dans $E$ est un \uline{homéomorphisme} de $E$ sur lui-même, car c'est une bijection continue dont la bijection réciproque $x\longrightarrow s^{-1}.x$ est aussi continue.

Supposons en outre que $G$ opère \uline{transitivement} dans $E$, autrement dit que $E$ soit (au sens algébrique) un \uline{espace homogène} de $G$ (Alg., chap.I, §7,n\textsuperscript{o}6). Soient $x_{\unsure{o|0}}$ un point de $E$, $H_o$ le sous-groupe de $G$ formé des opérateurs laissant invariant $x_o$ ( en \uline{Alg}., Bourbaki n'a pas cru bon de donner un nom à ce sous-groupe ; faut-il rattraper celà, et \unsure{l}'appeler ``stabilisateur'' de $x_o$ ?). La surjection continue $s\longrightarrow s.x_o$ de $G$ sur $E$ se factorise canoniquement en
\diagram{\begin{tikzcd}[ampersand replacement=\&]
G\arrow[r,"f"]\&G/H_o\arrow[r,"g"]\&E
\end{tikzcd}}
où $f$ est l'application canonique de $G$ sur l'espace homogè\deleted{\unsure{p|g}}ne $G/H_o$ et $g$ la bijection $s.H_o\to s.x_o$ de $G/H_o$ sur $E$ ; on sait en outre (Chap.I, 2eme éd. §9, n\textsuperscript{o}2, th.1) que si on munit $G/H_o$ de la topologie quotient de celle de $G$ par la relation d'équivalence $s^{-1}t\in H_o$, $g$ est \uline{continue}. On dit que $E$ est un \uline{espace homogène topologique}(correspondant
% archive sheet number: 95

\scanpage{97}

$(X_i)_{i\in\unsure{I|J}}$ système projectif filtrant de préschémas quasi compacts quasi séparés, $\exists i$ tq. $j\geq i\Rightarrow u_{ij}$ est affine.

Donc $X=\varprojlim X_i$ existe.

\deleted{\uline{Je dis que}}\inserted{A-t-on} $\operatorname{Cons}X\xrightarrow{\sim}\varprojlim\operatorname{Cons}(X_i)$ ???

En effet, soit $T$ un préschéma quasi compact et quasi-séparé, \unsure{alors}\ \deleted{\unsure{nous}} on a

{\scriptsize
\[
\begin{aligned}
\operatorname{Hom}(\operatorname{Cons}T,\operatorname{Cons}X)
&=\operatorname{App}\operatorname{cons}(T,X)=
\varinjlim_{\substack{\unsure{\mathcal P|\mathscr P|\mathfrak P}\ \text{une}\\\text{\unsure{une|{}} partition finie de}\\T\ \text{en ens loc.}\\\text{fermés constructibles}}}
\operatorname{Hom}\bigl(\deleted{\unsure{T_1|T_i|\Pi}}\coprod_{\unsure{\alpha|i}\in\mathrm{Ind}(\mathcal P)}T_i^{\mathcal P},X\bigr)\\
&=\varinjlim_{\mathcal P}\varprojlim_i\prod_{\alpha\in\mathrm{Ind}(\mathcal P}
\operatorname{Hom}(T_\alpha^{\mathcal P},X_i)
\end{aligned}
\]
\[
\operatorname{Hom}\bigl(\operatorname{cons}(T),\deleted{\unsure{\operatorname{cons}X|\operatorname{cons}X_i|\operatorname{Cons}X}}\varprojlim_{\unsure{i|{}}}\operatorname{Cons}(X_i)\bigr)
=\varprojlim_i\varinjlim_{\mathcal P}\prod_{\alpha\in\mathrm{Ind}\,\mathcal P}\operatorname{Hom}(T_\alpha^{\mathcal P},\unsure{X_i|X'})
\]
}

Cependant, considérons un autre système comme $X$, \unsure{soit} $(Y_j)_{j\in J}$, d'où
\[
Y\unsure{:|=}\varprojlim Y_j
\]
quelles sont les \deleted{morphismes} apl. constructibles de $X$ dans $Y$ ? \deleted{Il faut d'abord} \inserted{\unsure{On a|Cons}}

\[
\operatorname{Hom}(\operatorname{Cons}X,\operatorname{Cons}Y)=
\varinjlim_{\substack{\mathcal P\ \text{partitions}\\\text{finies de}\\X\ \text{en parties loc.}\\\text{fermées constr.}\ X_\alpha^{\mathcal P}}}
\prod_\alpha\operatorname{Hom}\bigl(\deleted{\unsure{T|X}}X_\alpha^{\mathcal P},Y\bigr)
\]
% archive sheet number: 96

\scanpage{98}

\mspage{- 13 -}\hfill n\textsuperscript{o} 285

formée de sous-groupes distingués fermés de $G$, et qui\unsure{.} converge vers $e$ ; en outre, lorsque les conditions de la prop.24 sont satisfaites, $G_\alpha$ est isomorphe à $G/N_\alpha$. Inversement :

PROPOSITION 25 - \uline{Soient} $G$ \uline{un groupe topologique séparé}, $(H_\alpha)$ \uline{une base de filtre sur} $E$, \uline{formée de sous-groupes distingués fermés et convergente vers} $e$. \uline{Pour} $H_\beta\subset H_\alpha$, \uline{soit} $f_{\alpha\beta}$ \uline{l'application canonique de} $G/H_\beta$ \uline{sur} $G/H_\alpha$ ; \uline{alors} $(G/H_\alpha,f_{\alpha\beta})$ \uline{est un système projectif de groupes topologiques vérifiant les conditions de la prop.}24, \uline{et si} $g_\alpha$ \uline{est l'application canonique de} $G$ \uline{sur} $G/H_\alpha$, $g=(g_\alpha)$ \uline{est un isomorphisme de} $G$ \uline{sur un sous-groupe partout dense de} $G'=\varprojlim G/H_\alpha$.

Pour $H_\beta\subset H_\alpha$\unsure{;|,}(relation de préordre que l'on notera $\alpha\leq\beta$), $G/H_\alpha$ s'identifie canoniquement à $(G/H_\beta)/(H_\alpha/H_\beta)$ (prop.16), donc $f_{\alpha\beta}$ est un homomorphisme, et il est immédiat que si $\alpha\leq\beta\leq\gamma$, $f_{\alpha\beta}\circ f_{\beta\gamma}={}=f_{\alpha\gamma}$. En outre on a, pour $H_\beta\subset H_\alpha$, $g_\alpha=f_{\alpha\beta}\circ g_\beta$, donc $g$ est une représentation continue de $G$ dans $G'$ (prop.20), et comme l'intersection des $H_\alpha$ est réduite à $e$, $g$ est injective. Comme $g_\alpha$ est surjective, il en est de même de l'application canonique de $G'$ sur $G/H_\alpha$ ; la prop.24 montre en outre que $g(G)$ est partout dense dans $G'$. Enfin, soit $V$ un voisinage de $e$ dans $G$, et soit $\alpha$ un indice tel que $H_\alpha\subset V$ ; l'ensemble des $y'=g(y)$ tels que $g_\alpha(y)\in g_\alpha(V)$ est un voisinage $W$ de l'élément neutre dans le sous-groupe $g(G)$ de $G'$. Mais la relation $g_\alpha(y)\in g_\alpha(V)$ entraîne par définition $y\in VH_\alpha\subset V^2$ ; on voit donc que $g(V^2)$ est un voisinage de l'élément neutre dans $g(G)$, ce qui achève la démonstration.

On dit que $g$ est \uline{l'application canonique} de $G$ dans $G'\unsure{\underset{\cdot}{=}|=}\varprojlim G/H_\alpha$

On notera que l'on n'a pas nécessairement $g(G)=G'$ (exerc).
% archive sheet number: 97
